\documentclass[12pt]{article}
\usepackage[a3paper, margin=1in]{geometry}
\usepackage{amsmath, amssymb, amsthm, ulem, booktabs}

\begin{document}

\section*{Věta 2.17}

Necht' $\{a_n\}$ je posloupnost. Potom $\limsup a_n, \liminf a_n \in H\left(\left\{a_n\right\}\right)$
a platí $\forall y \in H\left(\left\{a_n\right\}\right) :\ \liminf a_n \leq y \leq \limsup a_n$.

\textbf{Důkaz:}

\begin{itemize}
  \item $\limsup a_n \in H\left(\left\{a_n\right\}\right)$:
    \begin{itemize}
      \item $\limsup a_n = \infty$:
        \begin{quote}
          Dokažme konstrukcí vybrané posloupnosti $\{a_{n_k}\}$ s limitou $\infty$. Z definice víme, že $\{a_n\}$ není shora omezená.
          Sestavme indukcí $\{a_{n_k}\}$ t.ž. $a_{n_k} > k$.
          
          Protože posloupnost není omezená, nalezněme $n_1 \in \mathbb{N}$ takové, že $a_{n_1} > 1$.

          Uvažujme již vybraný index $n_{k-1}$, platí, že množina $\{a_n | n>n_{k-1}\}$ je stále shora neomezená.
          Pak zvolme $n_k > n_{k-1}$ splňující $a_{n_k} > k$.

          Máme posloupnost $a_{n_k} > k$ a triviálně z věty o dvou strážnících plyne $\lim_{k\to\infty} a_{n_k} = \infty$.
        \end{quote}
      \item $\limsup a_n = A \in \mathbb{R}$:
        \begin{quote}
          Dokažme konstrukcí vybrané posloupnosti $\{a_{n_k}\}$ s limitou $A$.
          Z definice platí $\limsup a_n = \lim_{n\to\infty} sup\{a_k | k\geq n\}= A$.
          Zaved'me posloupnost $S_n = sup\{a_k | k \geq n\}$.

          Označme lemma (*): $\forall \epsilon \in \mathbb{R}^+ \ \forall n \in \mathbb{N} \ \exists m \in \mathbb{N}, m>n:\ A-\epsilon < a_m < A+\epsilon$
          \begin{quote}
            \textbf{Důkaz:} Zvolme libovolný $\epsilon \in \mathbb{R}^+$. $S_n$ je nerostoucí, pak nalezněme $n_0 \in \mathbb{N}$ t.ž.
              \[ \forall n \in \mathbb{N}, n\geq n_0:\ A-\epsilon < A \leq S_n < A+\epsilon \]
              Zvolme libovolné $n\in\mathbb{N}, n\geq n_0$ a položme $n_* = n+1$.
              Z definice suprema a $S_n$ víme, že zároveň musí existovat $m\in\mathbb{N}, m\geq n_*$ takové, že
              \[ A-\epsilon \leq S_{n_*} - \epsilon < a_m \leq S_{n_*}\]
              \[ \Rightarrow A-\epsilon < a_m < A+\epsilon\ \land \ m\geq n_* > n\]
              Pro $n\in\mathbb{N}, n < n_0$ triviálně položme $n_* = n_0$ a podle předchozí úvahy opět existuje takové $m\in\mathbb{N}$.
              Epsilon i $n$ bylo libovolně volené, z čehož dostáváme
              $\forall\epsilon\in\mathbb{R}^+\ \forall n\in\mathbb{N}\ \exists m\in\mathbb{N}, m>n:\ A-\epsilon < a_m < A+\epsilon$.
          \end{quote}

          Konstruujme indukčně posloupnost $\{a_{n_k}\}$:
          \begin{itemize}
            \item Položme $\epsilon = 1, n=1$ a podle (*) nalezněme $n_1\in\mathbb{N}, n_1>1$ takové, že $A-1 < a_{n_1} < A+1$.
            \item Pro $n_k$ předpokládejme, že máme již zvolené $n_{k-1}$. Položme $\epsilon = \frac{1}{k}, n = n_{k-1}$.
              Podle (*) nalezněme $n_k\in\mathbb{N}, n_k>n_{k-1}$ splňující $A-\frac{1}{k} < a_{n_k} < A+\frac{1}{k}$.
          \end{itemize}

          Pak pro libovolně zvolený $\epsilon \in \mathbb{R}^+$ nalezněme $k_0 > \frac{1}{\epsilon}$,
          pak platí:
          \[ \forall k\in\mathbb{N}, k\geq k_0:\ |a_{n_k} - A| < \frac{1}{k} \leq \frac{1}{k_0} < \epsilon \]

          A tedy z definice platí, že limita posloupnosti $\{a_{n_k}\}$ je rovna $A$.
        \end{quote}
    \end{itemize}
  \item $\liminf a_n \in H\left(\left\{a_n\right\}\right)$: obdobně
  \item $\forall y \in H\left(\left\{a_n\right\}\right) :\ \limsup a_n \geq y$:
    \begin{itemize}
      \item $\limsup a_n = \infty$: Plyne triviálně z vlastnosti nekonečna.
      \item $\limsup a_n = A \in \mathbb{R}$:
        \begin{quote}
          Pro spor uvažujme $y > \limsup a_n$. Pak existuje vybraná posloupnost $\{a_{n_k}\}_{k=1}^\infty$
          s limitou $y$. Položme $\epsilon = \frac{y-A}{2}$ a nalezněme $n_1,n_2 \in\mathbb{N}$ takové, že:
          \[ \forall k\in\mathbb{N}, k\geq n_1:\ y- \frac{y-A}{2} < a_{n_k} < y+ \frac{y-A}{2} \ \Rightarrow \ \frac{y+A}{2} < a_{n_k} < \frac{3y-A}{2} \]
          \[ \forall n\in\mathbb{N}, n\geq n_2:\ A- \frac{y-A}{2} < S_n < A+ \frac{y-A}{2} \ \Rightarrow \ \frac{3A-y}{2} < S_n < \frac{y+A}{2} \]
          Položme $n_0 = max\{n_1,n_2\}$, pak platí:

          \[ \forall k\in\mathbb{N}, k\geq n_0:\ a_{n_k} > \frac{y+A}{2} > S_{n_k}\]
          Ale $S_{n_k}$ je supremum, čili toto je spor, jelikož musí platit $a_{n_k} \leq S_{n_k}$.
        \end{quote} 
    \end{itemize}
  \item $\forall y \in H\left(\left\{a_n\right\}\right) :\ \liminf a_n \leq y$: obdobně $\quad \square$
\end{itemize}

\end{document}
